\documentclass{article}
\usepackage{mathtools} % loads amsmath
\usepackage{amssymb}
\usepackage{amsthm}
\usepackage{hyperref}
\usepackage{cleveref} % must come after hyperref

\newtheorem{theorem}{Theorem}
\newtheorem{lemma}[theorem]{Lemma}
\theoremstyle{definition}
\newtheorem{definition}[theorem]{Definition}

\title{Project Euler Problem 6: Sum Square Difference}
\author{Brinhasavlin}
\date{9/30/2026}

\begin{document}
\maketitle


\section*{Problem}

The sum of the squares of the first ten natural numbers is,
\[
  1^2 + 2^2 + \cdots + 10^2 = 385.
\]
The square of the sum of the first ten natural numbers is,
\[
{(1 + 2 + \cdots + 10)}^2 = 55^2 = 3025.
\]
Hence, the difference between the sum of the squares of the first ten natural
numbers and the square of the sum is $3025 - 385 = 2640$.

Find the difference between the sum of the squares of the first one hundred
natural numbers and the square of the sum.


\section*{Proof}
This solution will derive from two known properties. There is a closed form for both the sum of squared integers and the squared sum of integers.

\begin{lemma}\label{lem:sum}
  For all $n \in \mathbb{N}$,
  \[
    \sum_{k=1}^{n} k = \frac{n(n+1)}{2}.
  \]
\end{lemma}
\begin{proof}
  We proceed by induction on $n$.

  \textbf{Base case} ($n = 1$).
  The left-hand side is $\sum_{k=1}^{1} k = 1$, and the right-hand side is
  \[
    \frac{1(1+1)}{2} = \frac{2}{2} = 1,
  \]
  so the statement holds for $n = 1$.

  \textbf{Inductive step.}
  Suppose that for some $n \ge 1$,
  \begin{equation}\label{eq:ih-sum}
    \sum_{k=1}^{n} k = \frac{n(n+1)}{2}. \tag{IH}
  \end{equation}
  We show the statement holds for $n+1$:
  \begin{align*}
    \sum_{k=1}^{n+1} k
      &= \sum_{k=1}^{n} k + (n+1)  && \text{(split off last term)} \\
      &= \frac{n(n+1)}{2} + (n+1)  && \text{by~\eqref{eq:ih-sum}} \\
      &= \frac{n(n+1) + 2(n+1)}{2} \\
      &= \frac{(n+1)(n+2)}{2}.
  \end{align*}
  Hence, the statement holds for all $n \in \mathbb{N}$.
\end{proof}

\begin{lemma}\label{lem:sumsq}
  For all $n \in \mathbb{N}$,
  \[
    \sum_{k=1}^{n} k^2 = \frac{n(n+1)(2n+1)}{6}.
  \]
\end{lemma}
\begin{proof}
  We proceed by induction on $n$.

  \textbf{Base case} ($n = 1$).
  The left-hand side is $\sum_{k=1}^{1} k^2 = 1^2 = 1$, and the right-hand side is
  \[
    \frac{1(1+1)(2 \cdot 1+1)}{6} = \frac{6}{6} = 1,
  \]
  so the statement holds for $n = 1$.

  \textbf{Inductive step.}
  Suppose that for some $n \ge 1$,
  \begin{equation}\label{eq:ih-square}
    \sum_{k=1}^{n} k^2 = \frac{n(n+1)(2n+1)}{6}. \tag{IH}
  \end{equation}
  We show the statement holds for $n+1$:
  \begin{align*}
    \sum_{k=1}^{n+1} k^2
      &= \sum_{k=1}^{n} k^2 + {(n+1)}^2      && \text{(split off last term)} \\
      &= \frac{n(n+1)(2n+1)}{6} + {(n+1)}^2  && \text{by~\eqref{eq:ih-square}} \\
      &= \frac{n(n+1)(2n+1) + 6{(n+1)}^2}{6} \\
      &= \frac{(n+1)\bigl(n(2n+1) + 6(n+1)\bigr)}{6} \\
      &= \frac{(n+1)(2n^2 + 7n + 6)}{6} \\
      &= \frac{(n+1)(n+2)(2n+3)}{6}.
  \end{align*}
  Hence, the statement holds for all $n \in \mathbb{N}$.
\end{proof}

\begin{theorem}
  For all $n \in \mathbb{N}$,
  \[
    {\left(\sum_{k=1}^{n} k\right)}^2 - \sum_{k=1}^{n} k^2 = \frac{n(n^2-1)(3n+2)}{12}.
  \]
\end{theorem}
\begin{proof}
  By \cref{lem:sum,lem:sumsq},
  \begin{align*}
    {\left(\sum_{k=1}^{n} k\right)}^2 - \sum_{k=1}^{n} k^2
      &= {\left(\frac{n(n+1)}{2}\right)}^2 - \frac{n(n+1)(2n+1)}{6} \\
      &= \frac{3{\bigl(n(n+1)\bigr)}^2 - 2n(n+1)(2n+1)}{12} \\
      &= \frac{n(n+1)\bigl(3n(n+1) - 2(2n+1)\bigr)}{12} \\
      &= \frac{n(n+1)(3n^2 - n - 2)}{12} \\
      &= \frac{n(n+1)(n-1)(3n+2)}{12} \\
      &= \frac{n(n^2-1)(3n+2)}{12}.
  \end{align*}
\end{proof}

\section*{Answer}
Applying the theorem with $n = 100$,
\begin{align*}
  {\left(\sum_{k=1}^{100} k\right)}^2 - \sum_{k=1}^{100} k^2
    &= \frac{100(100^2 - 1)(3 \cdot 100 + 2)}{12} \\
    &= \frac{100 \cdot 9999 \cdot 302}{12} \\
    &= \frac{301969800}{12} \\
    &= 25164150.
\end{align*}

\end{document}
